\documentclass[10pt]{article}
\usepackage[margin=1in]{geometry}
\usepackage{hyperref}
\usepackage{booktabs}
\usepackage{tabularx}
\usepackage{amsmath}
\usepackage{array}
\usepackage{url}
\usepackage{float}

\begin{document}

\title{Closed-Loop Threat-Guided\\Auto-Fixing of Kubernetes\\Container Security Misconfigurations\\Supplemental Appendices}
\date{}
\maketitle

\sloppy

\appendix
\section*{How to Read This Appendix (Plain English)}
This appendix keeps the explanations simple and points to the exact files that back every claim. Tables sit directly under their section titles so you can skim them quickly.

\section{Grok/xAI Failure Analysis}
\label{app:grok_failures}
The raw data for the Grok/xAI failure analysis can be found in \href{https://github.com/bmendonca3/k8s-auto-fix/blob/tcc-2025-12-0666-packet-2026-05-31/data/grok_failure_analysis.csv}{\texttt{data/grok\_failure\_analysis.csv}}. This file provides a comprehensive list of all failure causes and their corresponding counts, generated from the analysis of the 5,000-manifest Grok corpus.

\section{Risk Score Worked Example}
\label{app:risk_example}
The released telemetry enables reviewers to recompute risk units and $\Delta R/t$ for any queue item. As a concrete example we trace detection \texttt{001} from the Grok/xAI replay:
\begin{enumerate}
    \item Look up the detection metadata in \href{https://github.com/bmendonca3/k8s-auto-fix/blob/tcc-2025-12-0666-packet-2026-05-31/data/batch_runs/detections_grok200.json}{\texttt{data/batch\_runs/detections\_grok200.json}} to confirm the violation is \texttt{latest-tag}.
    \item Normalise the policy identifier and pull its risk weight and expected latency from \href{https://github.com/bmendonca3/k8s-auto-fix/blob/tcc-2025-12-0666-packet-2026-05-31/data/policy_metrics_grok200.json}{\texttt{data/policy\_metrics\_grok200.json}}. For \texttt{no\_latest\_tag} the risk is 50 units and the proposer+verifier expected time is 9.363~s (averaged from the recorded latencies).
    \item Inspect the proposer/verifier records (\href{https://github.com/bmendonca3/k8s-auto-fix/blob/tcc-2025-12-0666-packet-2026-05-31/data/batch_runs/patches_grok200.json}{\texttt{data/batch\_runs/patches\_grok200.json}}; \href{https://github.com/bmendonca3/k8s-auto-fix/blob/tcc-2025-12-0666-packet-2026-05-31/data/batch_runs/verified_grok200.json}{\texttt{data/batch\_runs/verified\_grok200.json}}) to see that the patch was accepted with a measured end-to-end latency of 7.339~s and verifier latency of 0.332~s.
\end{enumerate}
Because the patch succeeded, the pre-risk $R^{\text{pre}}=50$ drops to $R^{\text{post}} = 0$, yielding $\Delta R = 50$ and $\Delta R/t = 50 / 9.363 = 5.34$ risk units per second. Summing the same quantities across the corpus reproduces the risk-calibration table in the manuscript, as computed by \href{https://github.com/bmendonca3/k8s-auto-fix/blob/tcc-2025-12-0666-packet-2026-05-31/scripts/risk_calibration.py}{\texttt{scripts/risk\_calibration.py}}.

\noindent\textbf{Plain version.} We look up the finding, note its “risk points” and expected time to fix, check that the patch actually worked, and then see that the risk points went from 50 down to 0. That is the same math used for every item in the tables.

\section{Acronym Glossary}
\label{app:acronyms}
\begin{table}[H]
\centering
\small
\begin{tabularx}{\linewidth}{@{}l X@{}}
\toprule
\textbf{Acronym} & \textbf{Definition} \\
\midrule
CIS & Center for Internet Security \\
PSS & Pod Security Standards \\
CRD & Custom Resource Definition \\
RBAC & Role-Based Access Control \\
CTI & Cyber Threat Intelligence \\
KEV & CISA Known Exploited Vulnerabilities \\
EPSS & Exploit Prediction Scoring System \\
CVE/CVSS & Common Vulnerabilities and Exposures / Scoring System \\
RAG & Retrieval-Augmented Generation \\
MTTR & Mean Time To Remediate \\
CEL & Common Expression Language (Kubernetes) \\
SAST & Static Application Security Testing \\
DAST & Dynamic Application Security Testing \\
\bottomrule
\end{tabularx}
\end{table}

\section{Artifact Index}
\label{app:artifact_index}
\begin{table}[H]
\centering
\small
\caption{Primary artifacts bundled with the paper.}
\label{tab:primary_artifacts}
\begin{tabularx}{\linewidth}{@{}>{\ttfamily\raggedright\arraybackslash}p{0.58\linewidth} X@{}}
\toprule
\textbf{Artifact Path} & \textbf{Description} \\
\midrule
\texttt{data/live\_cluster/results\_1k.json} & Live-cluster replay outcomes (1,000 manifests, dry-run/live apply parity). \\
\texttt{data/batch\_runs/grok\_5k/metrics\_grok5k.json} & Grok/xAI telemetry (acceptance, latency, token counts) for the 5k sweep. \\
\texttt{data/risk/risk\_calibration.csv} & Risk accounting summary ($\Delta R$, residual risk, $\Delta R/t$) for supported and 5k corpora. \\
\texttt{data/metrics\_schedule\_compare.json} & Queue replay statistics for FIFO vs. risk-aware schedulers (rank, wait, $\Delta R/t$). \\
\texttt{data/grok\_failure\_analysis.csv} & Grok failure taxonomy (dry-run retrievals, StatefulSet validation, etc.). \\
\bottomrule
\end{tabularx}
\end{table}

\section{Evaluation Artifact Manifest}
\label{app:artifact_manifest}
\noindent Table~\ref{tab:artifact_manifest} expands the artifact index with record counts and purposes for full reproducibility.
\begin{table}[H]
\centering
\small
\caption{Key evaluation artifacts with record counts and purposes for full reproducibility.}
\label{tab:artifact_manifest}
\begin{tabularx}{\linewidth}{@{}>{\ttfamily\raggedright\arraybackslash}p{0.58\linewidth} >{\raggedright\arraybackslash}p{0.30\linewidth} r@{}}
\toprule
\textbf{Artifact Path} & \textbf{Purpose} & \textbf{Count} \\
\midrule
\texttt{data/live\_cluster/results\_1k.json} & Live-cluster replay outcomes (dry-run + apply) & 1{,}000 \\
\texttt{data/live\_cluster/summary\_1k.csv} & Live-cluster aggregate statistics & 1 \\
\texttt{data/batch\_runs/grok\_5k/metrics\_grok5k.json} & Grok-5k acceptance \& token telemetry & 5{,}000 \\
\texttt{data/batch\_runs/grok\_full/metrics\_grok\_full.json} & Manifest slice (1{,}313) acceptance & 1{,}313 \\
\texttt{data/batch\_runs/grok200\_latency\_summary.csv} & Proposer latency summary (Grok-200) & 280 \\
\texttt{data/batch\_runs/verified\_grok200\_latency\_summary.csv} & Verifier latency summary (Grok-200) & 140 \\
\texttt{data/eval/significance\_tests.json} & Statistical significance tests (z-test; Mann--Whitney~U) & 12 \\
\texttt{data/eval/table4\_counts.csv} & Table 4 manifest counts per corpus & 4 \\
\texttt{data/eval/table4\_with\_ci.csv} & Wilson 95\% confidence intervals & 4 \\
\texttt{data/scheduler/fairness\_metrics.json} & Scheduler fairness (Gini, starvation) & 830 \\
\texttt{data/scheduler/metrics\_schedule\_sweep.json} & Scheduler parameter sweep results & 16 \\
\texttt{data/risk/risk\_calibration.csv} & Risk reduction ($\Delta R$) per corpus & 2 \\
\bottomrule
\end{tabularx}
\end{table}

\section{Corpus Mining and Integrity}
\label{app:corpus}
\noindent\textbf{What this section is.} How we gathered the manifests and proved the dataset is intact.

\noindent\textbf{ArtifactHub mining pipeline.} Running the data collection helper\footnote{Command: \texttt{python scripts/collect\_artifacthub.py --limit 5000}. Full source: \href{https://github.com/bmendonca3/k8s-auto-fix/blob/tcc-2025-12-0666-packet-2026-05-31/scripts/collect_artifacthub.py}{scripts/collect\_artifacthub.py}.} renders Helm charts directly from ArtifactHub using \texttt{helm\ template}, normalizes resource filenames, and writes structured manifests under \href{https://github.com/bmendonca3/k8s-auto-fix/blob/tcc-2025-12-0666-packet-2026-05-31/data/manifests/artifacthub/}{data/manifests/artifacthub/}. The script records fetch failures and chart metadata so regenerated datasets can be diffed against the published summary.

\medskip
\noindent\textbf{Corpus hashes.} After manifests are rendered, \href{https://github.com/bmendonca3/k8s-auto-fix/blob/tcc-2025-12-0666-packet-2026-05-31/scripts/generate_corpus_appendix.py}{\texttt{python scripts/generate\_corpus\_appendix.py}} emits \href{https://github.com/bmendonca3/k8s-auto-fix/blob/tcc-2025-12-0666-packet-2026-05-31/docs/appendix_corpus.md}{docs/appendix\_corpus.md}, a SHA-256 inventory of every manifest (including the curated smoke tests in \href{https://github.com/bmendonca3/k8s-auto-fix/blob/tcc-2025-12-0666-packet-2026-05-31/data/manifests/001.yaml}{data/manifests/001.yaml} and \href{https://github.com/bmendonca3/k8s-auto-fix/blob/tcc-2025-12-0666-packet-2026-05-31/data/manifests/002.yaml}{002.yaml}). This appendix enables reproducibility reviewers to verify corpus integrity and trace individual evaluation examples back to their Helm chart origins. In short: the hash list is a tamper check and a map back to each source.

\end{document}
